\documentclass[10pt,a4paper]{amsart}
\usepackage[margin=19mm]{geometry}
\usepackage{amsmath,amssymb,mathtools}
\usepackage[scr=rsfs,cal=euler]{mathalfa}
\usepackage{xfrac}
\usepackage[unicode,hidelinks,bookmarks=false]{hyperref}
\newcommand{\ie}{i.e.\ }
\newcommand{\cf}{cf.\ }
\newcommand{\resp}{respectively\ }
\newcommand{\Subsection}{\subsection}
\providecommand{\nobreakdash}{\nobreak\hbox{-}}
\setlength{\emergencystretch}{2em}
\multlinegap0pt
\usepackage{bm}
\usepackage{bbm}
\usepackage{dsfont}
\usepackage{xspace}
\usepackage{cancel}
\usepackage{mathtools}
\usepackage{cleveref}
\usepackage{amsthm}
\makeatletter
\@ifundefined{theorem}{\newtheorem{theorem}{Theorem}}{}
\@ifundefined{lemma}{\newtheorem{lemma}[theorem]{Lemma}}{}
\makeatother
\makeatletter
\expandafter\ifx\csname align\endcsname\relax
\renewenvironment{align}{\linenomathNonumbers\oldalign}{\oldendalign\endlinenomath}
\fi
\expandafter\ifx\csname equation\endcsname\relax
\renewenvironment{equation}{\linenomathNonumbers\oldequation}{\oldendequation\endlinenomath}
\fi
\makeatother
\title[Closedness of the singular locus and generation for derived ]{Closedness of the singular locus and generation for derived categories}
\author{Souvik Dey, Pat Lank}
\date{Source: arXiv:2403.19564v2 (CC BY 4.0)}
\begin{document}
\maketitle

\noindent\emph{Source and licence.} Closedness of the singular locus and generation for derived categories. Version arXiv:2403.19564v2 (CC BY 4.0), distributed under the Creative Commons Attribution 4.0 International licence, as recorded in the arXiv licence metadata for this exact version and shown on the abstract page https://arxiv.org/abs/2403.19564v2. Full rights notice: licenses/arxiv-2403.19564-CC-BY-4.0.txt.\par\smallskip

\noindent\textit{Editorial scope.} Counted unit: the Lemma whose source label is \texttt{lem:abelian\_thick\_to\_bounded\_derived\_thick} in the article (source0.tex lines 340--342, source label \texttt{lem:abelian\_thick\_to\_bounded\_derived\_thick}), delivered together with the article's own complete original proof of it (source0.tex lines 344--358, one \texttt{proof} environment of 15 lines). No mathematics is written, repaired or re-derived: statement and proof are the article's own text line for line. Required context retained uncounted: lem:thick\_abelian\_filtration (lines 322--324). Every definition, display and label of those lines is retained unchanged. Omitted from this package and not redistributed: 1--321, 325--339, 343--343, 359--564. Nothing inside a retained excerpt is omitted or altered: every retained body is delivered exactly as the article's lines are written, so no omission receipt of any kind accompanies this package. Citations: the retained excerpts carry no citation command at all, so no bibliography entry of the article is transcribed and no reference is redistributed. Reference adaptation: none was needed and none was made; every reference inside the delivered unit resolves inside it, so no unresolved reference or citation is produced. Printed slips kept literal: no printed word, symbol, hypothesis or number of the counted statement or of its proof was altered. Layout and class: the article's own class is \texttt{amsart} with packages including inputenc, amssymb, amsmath, amsthm, enumerate, enumitem, colonequals, mlmodern, tikz-cd, microtype, stmaryrd, mathalpha, geometry, xcolor; the wrapper is the lane's pinned \texttt{amsart} editorial wrapper at 10pt on A4, which restates the theorem-like environments the retained text needs and adds the article's own macro definitions that the retained text needs (none), with extra wrapper packages (bm, bbm, dsfont, xspace, cancel, mathtools, cleveref). The delivered numbering is the unit's own. Assets: no retained span contains a figure environment, a graphics-inclusion command, a PDF-page inclusion, a table or a TikZ picture; no image file is copied into this package, the package declares no asset dependency and no image file is redistributed with it. Licence: version arXiv:2403.19564v2 of this article is distributed under the Creative Commons Attribution 4.0 International licence, as recorded in the arXiv licence metadata for this exact version and shown on the abstract page https://arxiv.org/abs/2403.19564v2. A full rights notice is delivered at licenses/arxiv-2403.19564-CC-BY-4.0.txt. Characters: every retained excerpt is pure ASCII, so the delivered engine, XeLaTeX with its default Latin Modern text font, reports no missing character.
\begin{lemma}\label{lem:thick_abelian_filtration}
    For any subcategory $\mathcal{S}$ of $\mathcal{A}$, we have $\operatorname{thick}(\mathcal{S})= \bigcup^\infty_{i=0} \operatorname{th}^i (\mathcal{S})$.
\end{lemma}

\begin{lemma}\label{lem:abelian_thick_to_bounded_derived_thick}
    If $\operatorname{thick}(\mathcal{S})=\mathcal{A}$, then $\langle \mathcal{S} \rangle = D^b(\mathcal{A})$.
\end{lemma}

\begin{proof}
    We know from \Cref{lem:thick_abelian_filtration} that $\operatorname{thick}(\mathcal{S})= \bigcup^\infty_{i=0} \operatorname{th}^i (\mathcal{S})$. Choose an object $E$ in $D^b (\mathcal{A})$. There exists a distinguished triangle in $D^b (\mathcal{A})$:
    \begin{displaymath}
        Z(E) \to E \to B(E)[1] \to Z(E)[1]
    \end{displaymath}
    where $Z(E)$ is the complex of cycles and $B(E)$ is the complex of boundaries. Note that the differentials of both $Z(E),B(E)[1]$ is zero, so these are complexes of shifts of objects in $\mathcal{A}$. If one can show that $Z(E),B(E)$ belongs to $\langle \mathcal{S} \rangle$, then the desired claim holds as one would have $E$ is in $\langle \mathcal{S} \rangle$. It suffices to check that each object of $A$ belongs to $\langle \mathcal{S} \rangle$. However, an induction argument tells us $\operatorname{th}^n (\mathcal{S})$ is contained in $\langle \mathcal{S} \rangle_{2^{n-1}}$ for all $n$. The case $n=1$ is clear, so assume there exists $n$ such that $\operatorname{th}^k (\mathcal{S})$ is contained in $\langle \mathcal{S} \rangle_{2^{k-1}}$ for all $1\leq k \leq n$. Let $E$ belong to $\operatorname{th}^{n+1} (\mathcal{S})$. There exists a short exact sequence in $\mathcal{A}$:
    \begin{displaymath}
        0 \to A \to B \to C \to 0
    \end{displaymath}
    where $E$ is a direct summand of one of the objects and the other two are objects of $\operatorname{th}^n (\mathcal{A})$. This gives us a distinguished triangle in $D^b (\mathcal{A})$:
    \begin{displaymath}
        A \to B \to C \to A[1].
    \end{displaymath}
    There are three cases, lets prove it for the case where $A,B$ are in $\operatorname{th}^n (\mathcal{S})$ as the others follow similarly. The inductive hypothesis tells us $A,B$ belong to $\langle \mathcal{S} \rangle_{2^{n-1}}$. Hence, it follows that $B$ belongs to $\langle \mathcal{S} \rangle_{2^n}$. This completes the proof.
\end{proof}

\end{document}
